\documentclass[11pt,a4paper]{article}
\usepackage[margin=1in]{geometry}
\usepackage[T1]{fontenc}
\usepackage[utf8]{inputenc}
\usepackage{graphicx}
\usepackage{hyperref}
\usepackage{enumitem}
\usepackage{float}
\usepackage{booktabs}
\usepackage{tabularx}

\title{Process Modeling and Process Mining\\Homework 2 — Process Mining Analysis}
\author{Peter Dobosi}
\date{\today}

\begin{document}
\maketitle

\section{Introduction and Goal}

This report presents a process mining analysis of two event logs from the BPI Challenge 2020 dataset, which records the travel expense reimbursement process at Eindhoven University of Technology (TU/e).

The central question driving this analysis is one posed by the challenge organizers themselves:

\begin{quote}
\textit{Is there a difference in throughput between domestic and international travel declarations?}
\end{quote}

To answer this, the analysis pursues the following objectives:
\begin{enumerate}[leftmargin=*]
    \item \textbf{Compare throughput} — measure and compare case durations (from submission to payment) between domestic and international declarations, both overall and at each approval step.
    \item \textbf{Identify bottlenecks} — determine where in the approval workflow the longest delays occur, and whether they differ between the two declaration types.
    \item \textbf{Discover and compare process models} — apply multiple discovery algorithms to both logs and evaluate the resulting models using conformance checking metrics.
\end{enumerate}

\section{Dataset Description}

\subsection{Source}

The dataset originates from the \textbf{BPI Challenge 2020}, the tenth International Business Process Intelligence Challenge organized by the IEEE Task Force on Process Mining. The data captures the travel expense reimbursement process at TU/e, covering 2017 (two departments) and 2018 (university-wide). It is publicly available at \url{https://data.4tu.nl/collections/BPI_Challenge_2020/5065541/1}.

\subsection{Process Overview}

At TU/e, domestic and international trips follow different procedures:
\begin{itemize}[leftmargin=*]
    \item \textbf{Domestic trips:} No prior permission is needed. The employee travels and submits a reimbursement claim afterwards.
    \item \textbf{International trips:} A travel permit must be approved by the supervisor before making any arrangements. Pre-paid travel costs can be claimed before the trip; the full declaration is submitted after the trip.
\end{itemize}

Both declaration types follow a similar approval workflow once submitted:
\begin{enumerate}[leftmargin=*]
    \item Submission by the employee.
    \item Approval by travel administration.
    \item Forwarding to the budget owner.
    \item Forwarding to the supervisor (skipped if the budget owner and supervisor are the same person).
    \item Optional director approval.
    \item Payment request and execution.
\end{enumerate}

At any approval step, a rejection can occur. The employee may then resubmit or withdraw the request.

\subsection{Event Logs Used}

Two of the five available logs were selected, corresponding to the two declaration types relevant to the throughput comparison:

\begin{table}[H]
\centering
\begin{tabular}{@{}lrrr@{}}
\toprule
\textbf{Log} & \textbf{Cases} & \textbf{Events} & \textbf{Events/Case (avg)} \\
\midrule
Domestic Declarations    & 10\,500 & 56\,437  & \textit{TODO} \\
International Declarations & 6\,449 & 72\,151 & \textit{TODO} \\
\bottomrule
\end{tabular}
\caption{Overview of the two event logs.}
\label{tab:logs}
\end{table}

The notably higher events-per-case ratio for international declarations (approximately \textit{TODO} vs.\ \textit{TODO}) already suggests a more complex process with more steps or more rework.

% TODO: After loading the data, fill in:
% - events/case averages
% - distinct activities per log
% - distinct resources per log
% - timespan per log
% - column names / attributes available

\section{Data Preparation}

% TODO: After loading the data, document:
% - Were there lifecycle transitions to filter?
% - Any missing values or anomalies?
% - Were any activities excluded or merged?
% - Any case-level filtering (e.g. by year)?
% - What formatting was applied (pm4py.format_dataframe)?

\textit{This section will be completed after loading and inspecting the raw event logs.}

\section{Exploratory Analysis}

\subsection{Activity Frequencies}

% TODO: For each log:
% - Bar chart of activity frequencies
% - Commentary on the most/least frequent activities
% - Include figures: output/domestic_activity_frequencies.png, output/international_activity_frequencies.png

\textit{TODO: Activity frequency analysis for both logs.}

\subsection{Case Duration Distribution}

% TODO: For each log:
% - Histogram of case durations
% - Mean, median, min, max
% - Compare the two distributions side by side
% - Include figures: output/domestic_case_durations.png, output/international_case_durations.png

\textit{TODO: Case duration distributions and comparison.}

\subsection{Process Variants}

% TODO: For each log:
% - Number of distinct variants
% - Coverage of top N variants
% - Include figures: output/domestic_variants.png, output/international_variants.png

\textit{TODO: Variant analysis for both logs.}

\subsection{Throughput Comparison}

This subsection addresses the central question directly.

% TODO:
% - Compare overall case durations (submission to payment) between domestic and international
% - Compare per-step durations using the performance DFG
% - Identify which steps contribute most to the difference
% - Include figures: output/throughput_comparison.png or similar

\textit{TODO: Direct throughput comparison between domestic and international declarations.}

\section{Process Discovery}

% TODO: For at least one of the logs (or both), run:
% - DFG (frequency and performance)
% - Alpha Miner
% - Heuristic Miner
% - Inductive Miner
% Visualize and discuss each model.
% Include figures for each.

\subsection{Directly-Follows Graph}

\textit{TODO: DFG discovery and visualization.}

\subsection{Alpha Miner}

\textit{TODO: Alpha Miner results.}

\subsection{Heuristic Miner}

\textit{TODO: Heuristic Miner results.}

\subsection{Inductive Miner}

\textit{TODO: Inductive Miner results.}

\section{Conformance Checking and Model Comparison}

% TODO:
% - Token-based replay: fitness, precision, generalization for each model
% - Comparison table
% - Bar chart (output/conformance_comparison.png)
% - Discussion: which model is best and why

\textit{TODO: Conformance checking results and model comparison.}

\section{Evaluation and Conclusion}

% TODO: Loop back to the three objectives from the introduction:
% 1. Throughput comparison: what did we find?
% 2. Bottlenecks: where are they, do they differ?
% 3. Model comparison: which algorithm worked best?
% Answer the central question: is there a throughput difference?

\textit{TODO: Final evaluation against stated goals.}

\end{document}
